Turning a 3D shape into a model made of interlocking toy bricks is a discrete design problem: every brick must sit on
studs, nothing may overlap, the result must hold together, and a person must be able to assemble it in order. We study
this problem under an explicit \emph{piece budget} and ask how fidelity scales with the number of bricks allowed. We
build a shared pipeline (voxelize, place, trim to one connected component, plan an assembly order, validate, and check
stability under a simplified static force model) and compare three generators inside it: a hand-written greedy search,
a 100k-parameter 3D convolutional policy trained by imitating that search, and the same policy fine-tuned with proximal
policy optimisation on the fidelity of the finished build. On a frozen benchmark of 53 held-out cases (procedural
shapes and Creative Commons Thingi10K meshes) at budgets of 50, 100 and 200 pieces, the learned pipeline reaches a mean
intersection-over-union of \meanRlVTwoNostopTwoHundred{} against the search's \meanGreedyVOneTwoHundred{} at the 200-piece cap, ties it on
procedural shapes, matches its validity and model-stability counts, and decodes in under 0.2\,s per shape; a second
training seed lands on the same number. Fidelity saturates between 100 and 200 pieces for every generator because
the cap rarely binds. Three negative results are reported in full: the learned stop action stays miscalibrated on
real meshes and is masked at decode time (a choice made on a held-out split, never on the benchmark); a stability term
in the reward is redundant with the hard constraints; and the connectivity inputs matter where the policy judges on
its own but not in the benchmark mean. All training ran in under five hours on one laptop. Code, benchmark records
and an audit that regenerates every number are released.
